\chap{The additive obstruction and the space of extensions}\label{sec:obstruction}
The strict definition gives a system of equations. In an additive setting, that system has an actual differential, and the failure to extend prescribed data has a quotient valued obstruction.

\subsection{The compatibility differential}
Assume each fiber \(A_v\) is an abelian group and each restriction is additive. Set
\[
 E=\prod_v A_v,\qquad R=\prod_{f:v\to w}A_v,
 \qquad (\delta x)_f=f^*(x_w)-x_v.
\]
The group \(R\) retains an entry for every arrow. It is not silently identified with a cohomology group.

\begin{proposition}[Differential identities]\label{KS08}
The map \(\delta:E\to R\) is additive, \(\Gamma(A)=\ker\delta\), and
\[
 (\delta x)_{1_v}=0,\qquad
 (\delta x)_{g\circ f}=f^*((\delta x)_g)+(\delta x)_f.
\]
\end{proposition}
\begin{proof}
Additivity follows coordinate by coordinate. Vanishing is exactly the defining coherence equation. For composition, add and subtract \(f^*x_w\):
\[
 f^*g^*x_z-x_v
 =f^*(g^*x_z-x_w)+(f^*x_w-x_v).
\]
The identity equation uses \((1_v)^*=\id\).
\end{proof}
\source{KS08}
This exact formula is more informative than a scalar error score. It specifies how defects transform and combine. Norms can be added later, but they do not replace the typed residual.

\subsection{The admissible corrections}
Fix \(U\subseteq\Ob\V\). Write \(p_U:E\to\prod_{v\in U}A_v\) for raw restriction. A correction preserving the prescribed values belongs to \(\ker p_U\). Its possible defects form the subgroup
\[
 B_U=\delta(\ker p_U)\le R.
\]
For a partial section \(s\), choose any raw extension \(\widetilde s\), for example zero off \(U\), and define
\begin{equation}\label{eq:obstruction}
 Q_U=R/B_U,\qquad
 \mathfrak o_U(s)=[\delta\widetilde s]\in Q_U.
\end{equation}
If another raw extension is chosen, their difference vanishes on \(U\), so its differential belongs to \(B_U\). Thus the class is independent of the lift.

\begin{theorem}[Extension criterion]\label{KS09}
For every coherent partial section,
\[
 \mathfrak o_U(s)=0\quad\Longleftrightarrow\quad\Ext_U(s)\ne\varnothing.
\]
\end{theorem}
\begin{proof}
If the class vanishes, \(\delta\widetilde s=\delta h\) for some \(h\in\ker p_U\). Then \(x=\widetilde s-h\) is coherent and restricts to \(s\). Conversely, if \(x\) is a coherent extension, \(h=\widetilde s-x\) vanishes on \(U\) and satisfies \(\delta h=\delta\widetilde s\). Hence the class is zero.
\end{proof}
\source{KS09}
The quotient is constructed before the solvability test. The theorem therefore does not define obstruction circularly by saying that it vanishes when a solution exists. It also does not give a decision procedure for membership in an arbitrary subgroup of an infinite product.

\plate{02_extension}{The raw family is repaired only along directions invisible on the prescribed views. The obstruction records whether its defect lies in the image of those admissible directions. Once an extension exists, a different kernel measures the remaining freedom.}{fig:extension}

\subsection{Existence is not uniqueness}
Define the ambiguity group
\[
 H_U=\ker\delta\cap\ker p_U.
\]
It consists of coherent changes that are invisible on the prescribed domain.

\begin{theorem}[Classification of a nonempty extension fiber]\label{KS10}
If \(x_0\in\Ext_U(s)\), translation gives an equivalence
\[
 H_U\simeq\Ext_U(s),\qquad h\longmapsto x_0+h.
\]
Consequently the extension is unique exactly when \(\mathfrak o_U(s)=0\) and \(H_U=0\).
\end{theorem}
\begin{proof}
Adding an element of \(H_U\) preserves coherence and the prescribed values. Conversely, the difference of two extensions is coherent and vanishes on \(U\). Subtraction of \(x_0\) is the inverse to translation. A nonempty translated group is a singleton exactly when the group is trivial.
\end{proof}
\source{KS10}
An extension fiber is therefore an affine copy, or torsor, of the invisible coherent group. Selecting one solution gives coordinates on the fiber. Without a selected solution, the affine structure does not designate an origin.

\begin{example}[Parity and an independent ambiguity]
For two views with \(f^*(m)=2m\) on \(\Z\), a prescribed value \(k\) at the coarse view extends exactly when \(k\) is even. The extension is \((k,k/2)\) and is unique. In particular \(1\) is obstructed and \(2\) extends nontrivially. The parity statement is an instance of Theorem~\ref{KS09} in the original source. If a separate, disconnected integer view is added and left unprescribed, every extendable seed has infinitely many completions. This last variant is a direct application of Theorem~\ref{KS10}, not a claim that the parity obstruction has changed.
\end{example}
\source{KS11}

\subsection{The exact sequence behind the criterion}
The additive structures make restriction a homomorphism
\(r_U:\Gamma(A)\to\Gamma_U(A)\). Using the zero extension, or lift independence, \(\mathfrak o_U\) is also a homomorphism.

\begin{proposition}[Exactness at partial sections]\label{KS12}
The sequence
\[
 \Gamma(A)\xrightarrow{r_U}\Gamma_U(A)
 \xrightarrow{\mathfrak o_U}Q_U
\]
is exact at \(\Gamma_U(A)\). Thus \(\ker\mathfrak o_U=\im r_U\), and restriction is onto precisely when \(\mathfrak o_U\) is the zero map.
\end{proposition}
\begin{proof}
A partial section is in the kernel precisely when its obstruction vanishes. By Theorem~\ref{KS09}, this means precisely that it is the restriction of a global section. The final equivalence is the statement that every partial section belongs to that image.
\end{proof}
\source{KS12}
No surjectivity onto the entire residual quotient is asserted. The residual group contains many entries that need not arise as the differential of a partial lift.

\subsection{Extension in stages}
For \(U\subseteq W\subseteq\Ob\V\), an extension first to \(W\) and then globally must retain the intermediate choice. There is an explicit equivalence
\begin{equation}\label{eq:stages}
 \Ext_U(s)\simeq
 \coprod_{\substack{t\in\Gamma_W(A)\\t|_U=s}}\Ext_W(t).
\end{equation}
It sends a global extension to its restriction to \(W\), together with itself as an extension of that restriction. The inverse forgets the intermediate coordinate. These operations are inverse because the required intermediate coordinate of a global section is its actual restriction. A coherent intermediate choice can still be obstructed, even when a different intermediate choice completes. Equation~\eqref{eq:stages} is why extension by stages is a dependent sum rather than a statement that every branch succeeds.
\source{KS13}

\chap{Protected correction as a general principle}\label{sec:protected}
The additive extension theorem is a particular instance of a useful general construction. Let \(E,R,S\) be abelian groups, with additive maps \(d:E\to R\) and \(p:E\to S\). The residual \(d\) measures compatibility, and \(p\) specifies the information that repair must preserve.

\begin{theorem}[Semantics preserving correction]\label{KS14}
For \(x\in E\), the following conditions are equivalent:
\[
\begin{aligned}
 &[dx]=0\text{ in }R/d(\ker p),\\
 &\exists h\in E:\ p(h)=0,\quad d(h)=d(x),\\
 &\exists y\in E:\ d(y)=0,\quad p(y)=p(x).
\end{aligned}
\]
The class of \(dx\) is unchanged when an admissible correction is subtracted from \(x\).
\end{theorem}
\begin{proof}
The first equivalence is quotient membership. Given \(h\), take \(y=x-h\). Given \(y\), take \(h=x-y\). If \(p(h)=0\), then \(d(x-h)-dx=-dh\) lies in the subgroup being quotiented out.
\end{proof}
\source{KS14}
Here \emph{semantics} means the declared mathematical observable \(p\). It is not an assertion that one homomorphism captures every meaning of every object. In a tower, \(p\) might preserve a seed operator and every new output component. In a graph, it might preserve designated vertex values. Different choices give different problems.

\begin{proposition}[A certificate of impossibility]\label{KS15}
Let \(L:R\to D\) be additive. If \(L\) vanishes on \(d(\ker p)\) but \(L(dx)\ne0\), no exact \(y\) can preserve \(p(x)\).
\end{proposition}
\begin{proof}
The map \(L\) descends to the obstruction quotient. Its nonzero value proves that the obstruction class is nonzero, so Theorem~\ref{KS14} excludes repair.
\end{proof}
\source{KS15}
This is often the most economical negative proof. One need not classify every possible correction. It is enough to find an invariant annihilating all admissible corrections but not the observed defect.

\subsection{Why unrestricted cancellation is empty information}
If every correction is permitted, taking \(h=x\) leaves the exact candidate zero. Thus \(\exists h:d(x-h)=0\) is automatic. It says nothing about recovering an object of interest. A repair theorem becomes informative only when it preserves something, is quantitatively controlled, or satisfies a comparison property that excludes arbitrary replacement.

\begin{example}[A rigid obstruction]
Take \(E=R=S=\Z\), \(d=p=\id\), and \(x=1\). The only admissible correction is zero, while the residual is nonzero. Hence there is no exact answer preserving the observed value. This is not a failure of the formalism. It is the correct certificate for an incompatible demand.
\end{example}

\subsection{What an obstruction does not yet provide}
An algebraic quotient does not automatically carry a useful norm. A zero infimum of distances to its boundary subgroup does not imply membership unless closure is controlled. Nor does vanishing supply a preferred correction when \(\ker d\cap\ker p\) is nontrivial. These two issues lead to different parts of the book: ambiguity is treated algebraically and through groupoids, while closed range and condition numbers enter the analytic theory.

\begin{insight}
A repair problem is not determined by its defect alone. It is determined by the defect \emph{and} the admissible changes. Keeping the pair \((d,p)\) explicit prevents a successful cancellation from being mistaken for a successful reconstruction.
\end{insight}
